(1) 盘和珠共享角速度，故 \(T=\frac12m\dot r^2+\frac12(I+mr^2)\dot\phi^2\)，弹簧势 \(V=k(r-r_0)^2/2\)。循环坐标 \(\phi\) 的广义动量为 \(p_\phi=(I+mr^2)\dot\phi=J\)，并由轴向无外力矩严格守恒。代入解得 \(\dot\phi(r)=J/(I+mr^2)\)。固定角动量约化可从 Routh 函数或总能量得到一维有效势，而不是把 \(\dot\phi\) 误当外部控制的常数。

(2) 径向能量写 \(E=m\dot r^2/2+U_{\rm eff}(r)\)，\(U_{\rm eff}(r)=k(r-r_0)^2/2+J^2/[2(I+mr^2)]\)。求导得 \(U'_{\rm eff}=k(r-r_0)-J^2mr/(I+mr^2)^2\)。平衡必须满足 \(k(r_e-r_0)=J^2mr_e/(I+mr_e^2)^2\)，故 \(r_e>r_0\)。将等式改写成 \(k[1-r_0/r_e]=J^2m/(I+mr_e^2)^2\)：左边对 \(r>r_0\) 严格递增、右边严格递减，且在两端分别一小一大，故唯一交点存在。

(3) 有效势二阶导数 \(U''_{\rm eff}=k-J^2m(I-3mr^2)/(I+mr^2)^3\)。在平衡处还可由 \(k-J^2m/(I+mr_e^2)^2=k r_0/r_e\) 改写为 \(U''_{\rm eff}(r_e)=kr_0/r_e+4J^2m^2r_e^2/(I+mr_e^2)^3>0\)。故径向小振动频率 \(\omega_r^2=U''_{\rm eff}(r_e)/m\)；角速度对小位移的线性响应 \(\delta\dot\phi=-2Jmr_e\delta r/(I+mr_e^2)^2\)，因此径向振动伴随旋转速率的反相微调。

(4) 系统总能量 \(E=\frac12m\dot r^2+\frac12k(r-r_0)^2+\frac{J^2}{2(I+mr^2)}\)，由于 \(U_{\rm eff}\) 在唯一平衡处有正曲率且当 \(r\to\infty\) 势能发散，该平衡是约束于固定 \(J\) 的稳定能量极小。若 \(I\to\infty\) 而 \(J\) 固定，盘角速度趋零，平衡 \(r_e\to r_0\)、\(\omega_r\to\sqrt{k/m}\)，退化为静止直槽中的普通弹簧滑块。若取极限时让 \(J/I\) 固定则得到另一个物理驱动条件，不能混同。
